% 「スーパーの店舗分布の二点相関関数」（notes/retail-cosmology.tex）の先の研究のノート。
% 扱うこと：九州の主な小売チェーンを並べて比べる（Issue #76）、店舗の三点相関関数（Issue #81）。
% 組み方：cd notes && latexmk retail-cosmology-beyond.tex（中間ファイルは notes/build/、ADR 006）
\documentclass[a4paper,11pt]{ltjsarticle}

\usepackage{amsmath,amssymb,mathtools,bm}
\usepackage{graphicx}
\usepackage{xcolor}
\usepackage[numbers,sort&compress]{natbib}
\usepackage{url}
\usepackage{booktabs}
\usepackage{hyperref}
\hypersetup{
  colorlinks=true,
  allcolors=blue!50!black,
  bookmarksnumbered=true,
  pdftitle={店舗分布の三点相関関数},
}

\graphicspath{{../docs/figures/}}
\renewcommand{\abstractname}{要旨}

\newcommand{\repourl}{https://github.com/inuverse44/retail-cosmology}
\definecolor{codecolor}{RGB}{0,95,135}
\newcommand{\code}[1]{\textcolor{codecolor}{\texttt{#1}}}

% 題は仮。九州のチェーンの比較も入るので、中身がそろったら決め直す。
\title{\textbf{店舗分布の三点相関関数}}
\author{Inuverse}
\date{\today}

\begin{document}
\maketitle

\begin{abstract}
準備中。
\end{abstract}

\clearpage
\tableofcontents
\clearpage

\section{序論}
\label{sec:intro}
% 前のノート（二点相関）の結果から始める：トライアルの 3–25 km の固まり、西友の数キロの固まり、
% Walmart の 5 km より近くの欠け、コスモスの固まりが九州の外から来ること。
% 二点相関は組の距離しか見ないので、固まりの形（丸い塊か、一列か、大型店のまわりの衛星か）は区別できない。
% 問い：(1) 三点相関で固まりの形がわかるか。(2) 九州の主なチェーンを並べると、並び方はどう違うか。
% 最後に構成の段落。

\section{枠組み}
\label{sec:framework}

\subsection{三点相関関数}
\label{sec:3pcf}
% 定義、三角形の取り方（r1, r2, θ）、reduced 三点相関 Q。
% 直感の図：丸い塊・一列・衛星の三つの並び方で、三角形の形の分布がどう変わるか。

\subsection{推定量}
\label{sec:3pcf-estimator}
% Szapudi--Szalay の推定量（DDD − 3DDR + 3DRR − RRR）/RRR。Landy--Szalay の三点版。
% θ について平均した低次の量（ルジャンドル展開）でノイズを抑える。

\subsection{多くの組を比べるときの判断}
\label{sec:global-envelope}
% n 社なら二社の組が n(n−1)/2。偶然「有意」に見える組を抑える。全体の包絡線検定。

\section{データ}
\label{sec:data}
% 九州の主な小売チェーン。店の数の目安（誤差 ∝ n^−0.6）と、どこまで言えるか。

\section{結果}
\label{sec:results}

\section{結論と議論}
\label{sec:conclusion}

% 最初の \cite を書いたらコメントを外す（文献は notes/references.bib を前のノートと共有する）。
% \bibliographystyle{unsrtnat}
% \bibliography{references}

\end{document}
